\section{Syntax}
\label{sec:grammar}

\subsection{Structure and modules}

A \ty{program} is a sequence of \ty{statement}s executed in order.
Source files may be combined with imports; an import splices the
imported file's text in front of the importing file's own code, under a
namespace alias:

\begin{lstlisting}
import io from "./std/io.y";   // io.f(...) addresses defs of io.y
\end{lstlisting}

\noindent An \ic{export \{ a, b \}} clause, when present, must be the
\emph{last} statement of a file. Paths resolve relative to the importing
file's directory.

\subsection{Declarations}

\begin{align*}
\ty{declaration} ::={}& \ic{let}~\ty{name}~[\ic{:}~\ty{type}]~\ic{=}~\ty{expr}~\ic{;} \\
  |~& \ic{const}~\ty{name}~[\ic{:}~\ty{type}]~\ic{=}~\ty{expr}~\ic{;} \\
  |~& \ic{def}~\ty{name}~\ic{(}~\ty{params}~\ic{)}~[\ic{->}~\ty{type}]~\ty{block}
\end{align*}

\ty{params} is a comma-separated list of \ty{name}~\ic{:}~\ty{type};
parameter annotations are \textbf{mandatory}, return types are optional
(inferred from the body, Section~\ref{sec:types}). A \ic{const}
declaration currently behaves exactly like \ic{let}; the keyword records
intent and is reserved for future immutability checking.

Function literals may also appear bound or anonymous:
\ic{let f = def (x: long) -> long \{ x \};}.

\subsection{Types}
\label{sec:type-grammar}

\begin{align*}
\ty{type} ::={}& \ic{long} \mid \ic{byte} \mid \ic{bool} \mid \ic{float}
              \mid \ic{str} \\
  |~& \ic{[}~\ty{type}~\ic{]}                    && \text{(heap array)}\\
  |~& \ic{[}~\ty{type}~\ic{;}~\ty{nat}~\ic{]}   && \text{(fixed array)}\\
  |~& \ty{type}~\ic{?}                           && \text{(option)}\\
  |~& \ic{*}~\ty{type}                           && \text{(pointer)}\\
  |~& \ic{(}~[\ty{type}~\{,\ty{type}\}]~\ic{)}~\ic{->}~\ty{type}
                                                 && \text{(function)}\\
  |~& \ic{\{}~[\ty{name}~\ic{:}~\ty{type}
        ~\{,~\ty{name}~\ic{:}~\ty{type}\}]~\ic{\}}
                                                 && \text{(record shape)}
\end{align*}

\noindent A record shape is structural; the shape \ic{\{ k: str \}} is
written to mean ``a record with at least a field \ic{k} of type
\ty{str}'' and is the conventional shape-hole annotation for AST-like
data. \ic{nat} is a non-negative integer literal. The type name
\ic{int} is accepted as a synonym of \ic{long} in annotations.

\subsection{Statements}

\begin{align*}
\ty{statement} ::={}& \ty{declaration} \\
  |~& \ty{expr}~\ic{;} \\
  |~& \ty{assignment}~\ic{;} \\
  |~& \ic{return}~[\ty{expr}]~\ic{;} \\
  |~& \ty{block} \\
  |~& \ic{if}~\ic{(}~\ty{expr}~\ic{)}~\ty{block}~[\ic{else}~(\ty{block} \mid \ty{statement})] \\
  |~& \ic{while}~\ic{(}~\ty{expr}~\ic{)}~\ty{block} \\
  |~& \ic{for}~\ic{(}~\ty{stmt}~\ic{;}~\ty{expr}~\ic{;}~\ty{stmt}~\ic{)}~\ty{block} \\
  |~& \ic{for}~\ty{name}~\ic{in}~\ty{expr}~\ic{..}~\ty{expr}~\ty{block} \\
  |~& \ic{foreach}~\ic{(}~\ty{name}~\ic{in}~\ty{expr}~\ic{)}~\ty{block} \\
  |~& \ty{block}
\end{align*}

\noindent \ty{block} is \ic{\{}~\ty{statement}\(^*\)~[\ty{expr}]~\ic{\}}:
a block is an \emph{expression} whose value is its trailing expression
(if any). A block must never be empty; write \ic{\{ 0; \}} where a
placeholder is needed. \ic{foreach} iterates the elements of an array;
the \ic{for} range form iterates \ic{lo..hi} with the low bound
inclusive and the high bound exclusive.

\subsection{Expressions}
\label{sec:expr-grammar}

\begin{align*}
\ty{primary} ::={}& \ty{long} \mid \ty{float} \mid \ty{string}
     \mid \ic{true} \mid \ic{false} \mid \ic{null} \mid \ty{name} \\
  |~& \ic{(}~\ty{expr}~\ic{)} \mid \ty{block} \\
  |~& \ic{[}~[\ty{expr}~\{,\ty{expr}\}]~\ic{]}
       \mid \ic{[}~\ty{expr}~\ic{;}~\ty{expr}~\ic{]} \\
  |~& \ic{\{}~[\ty{name}~\ic{:}~\ty{expr}~\{,\ldots\}]~\ic{\}}
       \quad \ic{\{}~[\ty{string}~\ic{:}~\ty{expr}~\{,\ldots\}]~\ic{\}} \\
  |~& \ic{def}~\ic{(}~\ty{params}~\ic{)}~[\ic{->}~\ty{type}]~\ty{block}
\end{align*}

\begin{align*}
\ty{postfix} ::={}& \ty{primary}
   ~\{~\ic{(}~[\ty{args}]~\ic{)} \mid \ic{[}~\ty{expr}~\ic{]}
   \mid \ic{[}~\ty{expr}\,?\ic{:}\,\ty{expr}\,?~\ic{]} \mid \ic{.}~\ty{name}
   \mid \ic{++} \mid \ic{--}~\} \\
\ty{unary} ::={}& [\ic{!} \mid \ic{-} \mid \ic{*} \mid \ic{\&}
   \mid \ic{++} \mid \ic{--}]~\ty{postfix}
\end{align*}

\subsection*{Operator precedence}

All binary operators are left-associative; lower rows bind tighter:

\begin{center}
\begin{tabular}{cll}
\toprule
level & operators & notes \\
\midrule
1 & \ic{||} & short-circuit \\
2 & \ic{\&\&} & short-circuit \\
3 & \ic{==} \ic{!=} & structural equality \\
4 & \ic{<} \ic{>} \ic{<=} \ic{>=} & numeric, or lexicographic \ty{str} \\
5 & \ic{::} & one-element cons (Section~\ref{sec:arrays}) \\
6 & \ic{+} \ic{-} & numeric, string/array \ic{+} \\
7 & \ic{*} \ic{/} \ic{\%} & numeric \\
\bottomrule
\end{tabular}
\end{center}

\subsection*{Assignment forms}

\begin{align*}
\ty{assignment} ::={}& \ty{place}~\ic{=}~\ty{expr} \\
  |~& \ty{place}~\ty{compound-op}~\ty{expr}
     && \ic{+=}~\ic{-=}~\ic{*=}~\ic{/=}~\ic{\%=} \\
  |~& \ty{place}~\ic{++} \mid \ic{++}~\ty{place}
   \mid \ty{place}~\ic{--} \mid \ic{--}~\ty{place}
\end{align*}

\noindent A \ty{place} is a name, an index expression \ic{a[e]} (arrays
and maps), or a field expression \ic{o.f} (records), possibly nested:
\ic{grid[0][1]}, \ic{rec.v[2]}, \ic{m['k'][i]}. Field access is
symbolic and exclusive: \ic{r.f} for records, \ic{m['k']} for maps.
Pointers add \ic{*p = e}.
Section~\ref{sec:exactly-once} defines the exact evaluation of every
form; all of them evaluate the place's subexpressions \textbf{exactly
once}.
